\section{Introduction}
\label{sec:intro}

World models – which capture an agent’s understanding of how the world changes with actions~\cite{ha2018world} – have emerged as a pivotal concept in robotics and Embodied AI. In embodied settings, a world model allows a robot to understand and predict its environment~\cite{ding2025understanding,long2025survey,fung2025embodied,li2025comprehensive}, and can function as the robot’s “internal brain”, enabling it to simulate future scenarios for planning and decision-making~\cite{chi2025wow,unifolmwma2025,wang2023anypose,Ko2023Learning}, or operate as an environment simulator~\cite{li2025vla,xiao2025world,liao2025genieenvisionerunifiedworld}.

Compared with cutting-edge spatial-prediction world models~\cite{genie3,team2025hunyuanworld}, embodied world models operate in context-rich environments that demand a deeper understanding of physical common sense. In robotics, the complexity and lack of standardization across setups further lead to a broad and diverse landscape of embodied-world-model designs. These models vary widely in their control conditions—ranging from approaches conditioned solely on images and language instructions~\cite{unifolmwma2025, chi2024eva, chi2025wow, chi2025mindlearningdualsystemworld}, to those incorporating keypoints~\cite{wang2025language}, trajectories~\cite{li2025manipdreamer3d, fu2025learning, jiang2025enerverse, quevedo2025evaluating, cen2025worldvla, guo2025ctrl}, depth, semantics, and other modalities~\cite{yang2025orv, li2025manipdreamer, zhou2024robodreamer}. They also differ in camera configurations, requiring either single-view inputs~\cite{chi2024eva, zhu2024irasim, cen2025worldvla, fu2025learning, unifolmwma2025, quevedo2025evaluating, wang2025language, chi2025mindlearningdualsystemworld, li2025manipdreamer3d} or multi-view setups~\cite{pmlr-v202-seo23a, chen2025robohorizonllmassistedmultiviewworld, pang2025learning, chi2025wow, liao2025genieenvisionerunifiedworld, yang2025orv, guo2025ctrl}. Moreover, several recent efforts have begun to pursue cross-embodiment generalization~\cite{zhi20253dflowaction, he2025scaling, chi2025wow}.
% 3D modeling requirements~\cite{}.
Therefore, despite the broader research of world models in general-purpose robotics, two core questions still remain: 
\begin{enumerate}
    \item Can these models generalize well enough to maintain perceptual fidelity from a human perspective?
    \item Are they robust and expressive enough to serve as universal priors for real-world embodied agents?
\end{enumerate}
% \textbf{\textit{\footnotesize (1)}} Can these models generalize well enough to maintain perceptual fidelity from a human perspective?\\
% \textbf{\textit{\footnotesize (2)}} Are they robust and expressive enough to serve as universal priors for real-world embodied agents?

Existing video-generation benchmarks~\cite{feng2024tc,ling2025vmbench,li2025worldmodelbench,yue2025ewmbench,liu2024evalcrafter,ji2024t2vbench,zheng2025vbench} largely target general-purpose settings or isolated dimensions and overlook the unique requirements of robotic world models. Most evaluations emphasize visual fidelity or coarse task success, but rarely assess deeper embodied abilities such as physical plausibility, planning rationality, and actionability. This gap makes progress difficult to measure: a model may score well on conventional video metrics~\cite{unterthiner2018towards,1284395,Oquab2023DINOv2,fardo2016formalevaluationpsnrquality} yet produce physically impossible or contextually incorrect predictions in robotic scenarios. Our results further confirm this misalignment—standard video-quality scores correlate poorly with human judgments in embodied settings—highlighting the need for more reliable evaluation standards.

Consequently, in this paper, we address these challenges by proposing a new comprehensive benchmark for embodied world models, and use it to systematically evaluate foundational models under the simplest \textit{image-to-video} setting. Our benchmark, \textbf{WoW-World-Eval}, as illustrated in Figure~\ref{fig:WoWbench_overall_design}, is designed around the core capabilities that an embodied world model should possess: perceiving the environment, understanding and planning based on task instructions, predicting and simulating future world states, executing real-world interactions, and generalizing across diverse scenarios and embodiments.
% understanding physical dynamics, predicting realistic action outcomes, maintaining semantic consistency with instructions, and enabling successful task execution in simulation. 
It contains about 609 robot manipulation samples with meticulous cleaning and annotation by human annotators. Also we incorporate 22 evaluation metrics aligned with our core dimensions across Video Quality, Instruction Understanding, Planning Reasoning, Physical Law, and Execution Accuracy.
% human perception and task performance – for instance, whether the predicted video obeys conservation of energy, or whether it leads to a correct task completion when interpreted by a control policy. 

% These metrics are  from the Human Turing Test: by collecting fine-grained human preference data on real vs. generated videos, we can calculate the proportion of videos generated by each model that can fool humans, and examine whether their growth trends correlate with the scores these models achieve on our benchmark. Therefore, on the other hand, we can certainly construct a Machine Turing Test, which in this case is Inverse Dynamic Model Turing Test. By determining whether videos generated by the model can “fool” the IDM, we test its real-world execution success rate. Since this IDM has only seen real-world execution videos, if the generated model can deceive the IDM into outputting actions that can be genuinely executed, it indicates that the generated videos are extremely similar to real ones.
Moreover, \textbf{WoW-World-Eval} follows the two-alternative forced-choice (2AFC)~\cite{fechner1860elemente} methodology from psychophysics to establish the evaluation as a standardized Turing Test for generative video models. By collecting fine-grained human answers distinguishing real and generated videos, we compute the proportion of generated videos from each model that successfully fool human evaluators. Notably, an overall human preference alignment score exceeding Pearson Correlation = 0.93 demonstrates the effectiveness of our benchmark in evaluating high-quality generations and serves as a reliable proxy for the \textit{Human Turing Test}.

In addition to the human-centered evaluation, we also introduce a \textit{Machine Turing Test}—specifically, an \textit{Inverse Dynamics Model (IDM) Turing Test}. In this setting, we assess whether videos generated by a model can “fool” an IDM that has only been trained on real-world execution sequences. If the generated videos lead the IDM to output plausible actions that are executable in the real world, it indicates that the model's outputs are indistinguishable from real data in terms of physical and action plausibility.

% This human-aligned metric serves as a stringent litmus test for world models’ realism. Models are thus judged not only by standard quality scores but also by whether they can fool human observers on our benchmark’s tasks. 
By evaluating existing models under this new benchmark, we reveal which models already exhibit credible world understanding and where they fall short. We believe this benchmark and the accompanying Turing Test criterion will provide a much-needed standard for the field, driving research towards embodied world models that can truly imagine the world with the accuracy and fidelity that robotics applications demand. Our contributions are threefold as follows:
\begin{itemize}
    \item A comprehensive World Model Benchmark, \textbf{WoW-World-Eval}, focuses on the Embodied AI domain, introducing a new perspective with a novel framework for the five core abilities in the embodied world model. 
    
    \item Based on \textbf{WoW-World-Eval}, we propose two novel Turing tests aligned with our benchmark, the Human Turing test and IDM Turing Test, that can distinguish the models' true ability as an embodied world model to simulate and interact with the real world.

    \item We curate 609 high-quality robot manipulation samples with careful human annotations in \textbf{WoW-World-Eval}. Using this benchmark, we conduct a comprehensive evaluation of various world models, whose performance provides new insights into the strengths, limitations, and generalization gaps of current embodied world models.
    % \item We are the first to test the capabilities of the video foundation model in real world using an inverse dynamic model, which better demonstrates its ability as a world model.
    % \item We introduce a sim-to-real simulator as our detector to measure how related that the World Model video is related to real-world execution(the sim-to-real simulator has already been verified to narrow the gap between the simulator and the real world) using a fixed IDM.
\end{itemize}

